\scanpage{38}

% B01 | printed letterhead, upper left
% letterhead
\begin{center}
UNIVERSITÉ DE PARIS\\
---$\cdot$---\\
ÉCOLE NORMALE SUPÉRIEURE\\
---$\cdot$---\\
45, RUE D'ULM\\
ODÉON 06$\cdot$45\\
---$\cdot$---
\end{center}

% B02 | Dém., first four manuscript lines
\uuline{Dém}. $G/S$ opère sans point fixe sur $E/S$ : si $gx=s_0x$, $g=s_0s$ avec $sx=x\Rightarrow g\in S$. Si $g\notin S$, $g$ transforme tout sommet de $E$ en un sommet distinct, donc tout simplexe en un simplexe disjoint. $E/G={}^{E/S}\!/\!{}_{G/S}$.

% B03 | short rule and underlined conclusion
--- \uline{Le Théorème 1} \deleted{\unsure{I|1}} résulte de ce lemme et de ce qui a été vu précédemment.

% B04 | short line below conclusion
\deleted{$\unsure{\text{que }\pi_1(B)|p_*\pi_1(B)|p_*\pi_1(E)}$}

% archive sheet number: 37

